\documentclass[a4paper,12pt]{article}
\usepackage{graphicx}
\usepackage{caption}
\usepackage{subcaption}
\usepackage{float}
\usepackage[margin=1in]{geometry}
\usepackage{lmodern}
\usepackage{amsmath}
\usepackage{xcolor}
\usepackage{listings}
\usepackage{pgffor}

\lstset{
  breaklines=true,
  postbreak=\mbox{\textcolor{red}{$\hookrightarrow$}\space},
}

\graphicspath{{../../python/kNN/k_optimization/results/}, {../../python/kNN/output/}}

\begin{document}

\title{Relatório kNN}
\author{Eduardo}
\date{\today}
\maketitle

\section*{Exercício 1}

As figuras a seguir apresentam a relação entre número de vizinhos, \texttt{k}, e a precisão média na validação /textit{k-fold} para três dados reais.

\begin{figure}[H]
  \centering
  \includegraphics[width=0.6\textwidth]{cdc_diabetes_health_indicators.png}
  \caption{Relação entre número de vizinhos (\texttt{k}) e precisão média para o conjunto de dados \texttt{cdc\_diabetes\_health\_indicators}.}
  \label{fig:cdc_diabetes_health_indicators}
\end{figure}

\begin{figure}[H]
  \centering
  \includegraphics[width=0.6\textwidth]{diabetic_retinopathy_debrecen.png}
  \caption{Relação entre número de vizinhos (\texttt{k}) e precisão média para o conjunto de dados \texttt{diabetic\_retinopathy\_debrecen}.}
  \label{fig:diabetic_retinopathy_debrecen}
\end{figure}

\begin{figure}[H]
  \centering
  \includegraphics[width=0.6\textwidth]{occupancy.png}
  \caption{Relação entre número de vizinhos (\texttt{k}) e precisão média para o conjunto de dados \texttt{occupancy}.}
  \label{fig:occupancy}
\end{figure}

Nota-se que o número ótimo de vizinhos muda conforme o conjunto de dados. Para o encontrar, uma forma eficiente e automática é a aplicação de um otimizador que busque \texttt{k} que maximize a precisão média.

\section*{Exercício 2}

As figuras a  seguir apresentam pares de superfície de decisão e espaço característico para diferentes conjuntos sinteticos.

\foreach \dataset in {blobs, moons, spirals, xor} {
  \begin{figure}[H]
    \centering
    \begin{subfigure}{0.45\textwidth}
      \centering
      \includegraphics[width=\textwidth]{\dataset\_decision\_boundary.png}
      \caption{Superfície de decisão para o conjunto de dados \texttt{\dataset}.}
    \end{subfigure}%
    \hfill
    \begin{subfigure}{0.45\textwidth}
      \centering
      \includegraphics[width=\textwidth]{\dataset\_characteristic\_space.png}
      \caption{Espaço característico para o conjunto de dados \texttt{\dataset}.}
    \end{subfigure}
    \caption{Superfície de decisão e espaço característico para o conjunto de dados \texttt{\dataset}.}
  \end{figure}
}

\end{document}
